%% \documentclass[../main.tex]{subfiles}
%% \usepackage{svg}
%% \graphicspath{{\subfix{../images/}}}
%% \begin{document}

%% TODO: Move common code outside of localized versions of the book.

\newcounter{example-counter}
\setcounter{example-counter}{1}

%% This procedure adds the "Example" block to the text.
\newcommand{\example}[1]{
  \vspace{8pt}
  \begin{tabularx}{\textwidth}{m{1cm} m{9cm}}
    \includesvg[width=1.25cm]{the-noun-project/request-mirrored}
    & \textbf{Example \arabic{example-counter}}: #1 \\
  \end{tabularx}
  \addtocounter{example-counter}{1}
}

\newcounter{experiment-counter}
\setcounter{experiment-counter}{1}

%% This procedure adds the "Experiment" block to the text.
\newcommand{\experiment}[2]{
  \vspace{8pt}
  \begin{tabularx}{\textwidth}{m{.15\textwidth}X m{.85\textwidth-4}X}
    \includesvg[width=1cm]{the-noun-project/flask}
    & \textbf{Experiment №\arabic{experiment-counter}:} #2 \\
  \end{tabularx}
  \addtocounter{experiment-counter}{1}
}

\newcommand{\note}[1]{
  \vspace{8pt}
  \begin{tabularx}{\textwidth}{m{1cm} m{9cm}}
    \includesvg[width=1cm]{the-noun-project/note}
    & \textbf{Note:} #1 \\
  \end{tabularx}
}

\newcommand{\hotkey}[1]{
  \texttt{#1}
}

%% This procedure allows to insert a music note.
%% Syntax:
%%   \musicnote{<octave>}{<note-name>}{<frequency>}
\newcommand{\musicnote}[3]{
  &
  \ifstrequal{#2}{C}{Do   & C#1}{}
  \ifstrequal{#2}{D}{Re   & D#1}{}
  \ifstrequal{#2}{E}{Mi   & E#1}{}
  \ifstrequal{#2}{F}{Fa   & F#1}{}
  \ifstrequal{#2}{G}{Sol  & G#1}{}
  \ifstrequal{#2}{A}{La   & A#1}{}
  \ifstrequal{#2}{B}{Ti   & B#1 (H#1)}{}
  & #3 \\
}

%% Taken from:
%%   <https://tex.stackexchange.com/questions/184923/how-to-include-a-second-file-only-if-environment-variable-is-set>
\newcommand{\newgetenv}[2][]{%
 \CatchFileEdef{\temp}{"|kpsewhich --var-value #2"}{\endlinechar=-1\relax}%
 \if\relax\detokenize{#1}\relax\temp\else\edef#1{\temp}\fi%
}%

\newcommand\esymbol[1]{%
  \begin{circuitikz}%
    \draw (0,0) to [#1] (1,0);%
  \end{circuitikz}%
}

\newcommand*{\soundWaveIcon}[0]{%
  \begin{tikzpicture}
    \draw[black, fill=black] (0, 0) circle (.25mm);
    \draw (.5mm, 0.7mm) arc (45:-45:1mm);
    \draw (1mm, 1mm) arc (45:-45:1.5mm);
    \draw (1.5mm, 1.3mm) arc (45:-45:2mm);
  \end{tikzpicture}%
}

%% Taken from:
%%   <https://tex.stackexchange.com/questions/40624/howto-typeset-the-direct-current-symbol-in-latex>
\newcommand{\textdirectcurrent}{%
  \settowidth{\dimen0}{$=$}%
  \vbox to .85ex {\offinterlineskip
    \hbox to \dimen0{\leaders\hrule\hfill}
    \vskip.35ex
    \hbox to \dimen0{%
      \leaders\hrule\hskip.2\dimen0\hfill
      \leaders\hrule\hskip.2\dimen0\hfill
      \leaders\hrule\hskip.2\dimen0
    }
    \vfill
  }%
}

%% This star is used to mark sections of the book that contain advanced
%% material.
\newcommand{\Star}{\ensuremath{^\star}\kern-\scriptspace}

%% \end{document}
